\section{Domain Analysis}
\subsection{Application Context and Utilized Technologies}

The primary objective of this laboratory work is to implement an \textbf{analog signal
acquisition and conditioning system} for a digital temperature sensor using
\textbf{FreeRTOS} on the ESP32 microcontroller. The application demonstrates a
complete sensor pipeline: periodic raw data acquisition, multi-stage signal
preprocessing (saturation, median filtering, weighted averaging), threshold-based
alert detection, and structured reporting on both an OLED display and the serial
interface.

The system uses the \textbf{DHT11} digital temperature and humidity sensor. The sensor
communicates over a proprietary single-wire protocol on GPIO\,4 and is read through the
\textit{Adafruit DHT} library. The DHT11 requires a minimum interval of approximately
2\,seconds between consecutive reads, which determines the acquisition task period.

The signal conditioning stage implements three classical preprocessing techniques
applied in sequence:
\begin{itemize}
  \item \textbf{Saturation (clamping):} The raw sensor value is clamped to the physical
    measurement range of the DHT11 ($0$--$50\,^\circ$C). Values outside this range are
    replaced with the nearest boundary, preventing physically impossible readings from
    propagating through the pipeline.
  \item \textbf{Median filter (salt-and-pepper removal):} A sliding window of size 5
    stores the most recent saturated samples. Each cycle, the window contents are sorted
    and the median value is selected. This non-linear filter is highly effective at
    removing isolated outlier spikes without distorting the underlying signal trend.
  \item \textbf{Exponential Moving Average (EMA):} A first-order IIR low-pass filter
    with coefficient $\alpha = 0.2$ is applied to the median-filtered output:
    \[
      y_n = \alpha \cdot x_n + (1 - \alpha) \cdot y_{n-1}
    \]
    This provides weighted smoothing where recent samples contribute more than older
    ones, reducing high-frequency noise while maintaining responsiveness.
\end{itemize}

After the full pipeline, alert thresholds with hysteresis (ON $\geq 28\,^\circ$C,
OFF $\leq 24\,^\circ$C) are evaluated on the median-filtered value to produce a
stable binary alert state.

Task synchronization is achieved through two FreeRTOS primitives:
\begin{itemize}
  \item A \textbf{binary semaphore} (\texttt{xSemNewData}) given by TaskAcquisition
    after each read and taken by TaskConditioning, ensuring the conditioning logic runs
    exactly once per new sample
  \item A \textbf{mutex} (\texttt{xSensorMutex}) protecting all shared sensor data
    fields against concurrent access by the three tasks
\end{itemize}

\subsection{Hardware and Software Components}

\begin{center}
\begin{tabular}{| p{4cm} | p{10cm} |}
\hline
\textbf{Component} & \textbf{Description \& Role} \\ \hline
\textbf{ESP32 DevKit} & Dual-core microcontroller running FreeRTOS. Hosts all three
  application tasks and provides GPIO and I2C peripherals. \\ \hline
\textbf{DHT11} & Digital temperature and humidity sensor on GPIO\,4. Single-wire
  protocol, $0$--$50\,^\circ$C range, $\pm 2\,^\circ$C accuracy, minimum 2\,s
  between reads. \\ \hline
\textbf{Green LED} & Connected to GPIO\,25. Lights up when temperature is within the
  normal range (alert OFF). \\ \hline
\textbf{Red LED} & Connected to GPIO\,26. Lights up when the temperature alert is active
  (temperature above $28\,^\circ$C). \\ \hline
\textbf{OLED 128$\times$64 SSD1306} & Connected on the I2C bus (SDA=GPIO\,21,
  SCL=GPIO\,22), address 0x3C. Primary display showing raw, filtered, and averaged
  values plus alert state at 500\,ms intervals. \\ \hline
\textbf{LCD 20$\times$4 I2C} & Connected on the same I2C bus, address 0x27. Backup
  display showing averaged temperature and alert state. \\ \hline
\textbf{$220\,\Omega$ Resistors} & Current-limiting resistors for each LED. \\ \hline
\textbf{$10\,\mathrm{k\Omega}$ Pull-up} & Pull-up resistor on the DHT11 data line
  to 3.3\,V for reliable communication. \\ \hline
\textbf{Serial-to-USB Interface} & UART at 115200 baud for serial plotter data
  (\texttt{>name:value} format) and periodic structured reports. \\ \hline
\textbf{FreeRTOS} & Built into the ESP32 Arduino core. Provides preemptive task
  scheduling, binary semaphore, and mutex. \\ \hline
\textbf{PlatformIO} & Build system and IDE integration used for compilation, upload,
  and serial monitoring. \\ \hline
\textbf{Serial Plotter (VS Code)} & VS Code extension for real-time plotting of data
  received via the serial port. Parses lines starting with \texttt{>} in
  \texttt{name:value} format and renders interactive time-series graphs, allowing
  side-by-side comparison of raw, median-filtered, and EMA-smoothed traces. \\ \hline
\textbf{Breadboard} & Prototyping base for connecting the DHT11, LEDs, OLED, LCD,
  and the ESP32 without soldering. \\ \hline
\end{tabular}
\end{center}

\subsection{System Architecture and Solution Justification}

The system is organized into three hardware-abstraction layers following the MCAL /
ECAL / SRV pattern:

\begin{enumerate}
  \item \textbf{MCAL (Microcontroller Abstraction Layer):} Contains only
    \texttt{GpioDriver}, which wraps \texttt{pinMode}, \texttt{digitalWrite},
    \texttt{digitalRead}, and \texttt{analogRead}. All direct register / Arduino
    hardware calls are isolated here.
  \item \textbf{ECAL (Electronic Component Abstraction Layer):} Contains component
    drivers: \texttt{TempSensor} (DHT11 via Adafruit DHT), \texttt{LedDriver},
    \texttt{OledDisplay} (Adafruit SSD1306), \texttt{LcdDisplay}
    (LiquidCrystal\_I2C), and \texttt{SerialIO}. Each module depends only on MCAL
    or its own library and exposes a clean C API. All pin assignments are declared
    in the respective header files for easy reference during circuit assembly.
  \item \textbf{SRV (Service Layer):} Contains \texttt{SensorData} (shared data store
    with mutex/semaphore, configurable pipeline parameters, and sensor catalogue),
    and the three FreeRTOS task implementations:
    \texttt{TaskAcquisition}, \texttt{TaskConditioning}, and \texttt{TaskReporter}.
\end{enumerate}

The data flow proceeds as follows:
\begin{enumerate}
  \item \textbf{TaskAcquisition} reads the DHT11 sensor every 2000\,ms, stores the
    raw temperature and validity flag in \texttt{SensorData} (under mutex), and gives
    the \texttt{xSemNewData} binary semaphore.
  \item \textbf{TaskConditioning} unblocks immediately on the semaphore, reads the
    latest raw temperature, and applies the three-stage pipeline: saturation $\to$
    median filter $\to$ EMA. It stores the intermediate and final values in
    \texttt{SensorData}, evaluates alert thresholds with hysteresis, and drives the
    green/red LEDs accordingly.
  \item \textbf{TaskReporter} runs every 500\,ms via \texttt{vTaskDelayUntil}, reads
    all shared values under mutex, updates the OLED (raw, median, EMA, alert) and LCD
    (backup), emits a serial plotter line, and every 5\,seconds prints a full
    structured report to STDIO.
\end{enumerate}

\subsection{Relevant Case Study}
A real-world application of multi-stage analog signal conditioning is
\textbf{environmental monitoring in greenhouse control systems}. Raw temperature
readings from sensors placed in a greenhouse are subject to several types of noise:
sudden spikes from direct sunlight on the sensor (salt-and-pepper noise), electrical
interference from nearby irrigation pumps, and gradual drift. A saturation stage
rejects physically impossible readings (e.g., $-40\,^\circ$C indoors), a median filter
removes the isolated spikes without smoothing away real temperature changes, and a
weighted moving average provides the stable baseline needed for climate control
decisions. Alert thresholds with hysteresis then trigger ventilation or heating
actuators without the rapid on/off cycling that would occur with a single-threshold
comparator -- a pipeline directly reflected in this laboratory work.
